\subsection{Absorption par porteurs libres}

Une fois les électrons promus dans la bande de conduction, ils se comportent comme des porteurs quasi-libres : le champ laser les accélère, et ils dérivent en opposition à $\vec{E}$~\cite{mao2004}. On pourrait s'attendre à ce que ces porteurs continuent à absorber des photons et à gagner de l'énergie directement. Cependant, une absorption simple-photon par un électron libre isolé est interdite, car elle ne peut pas satisfaire simultanément la conservation de l'énergie et de la quantité de mouvement.

Pour l'absorption d'un photon $(\omega_0,\vec{k}_{\text{ph}})$ par un électron de vecteur d'onde initial $\vec{k}_i$ vers un état final de vecteur d'onde $\vec{k}_f$, les lois de conservation de l'énergie et de la quantité de mouvement doivent être satisfaites. Elles s'écrivent :
\begin{equation}
  E_f = E_i + \hbar\omega_0 , \qquad \vec{k}_f = \vec{k}_i + \vec{k}_{\text{ph}}.
  \label{eq:const}
\end{equation}

En utilisant la bande parabolique $E = \hbar^2 k^2 / 2m^\ast$ et la quantité de mouvement du photon $k_{\text{ph}} = \omega_0/c = 2\pi/\lambda_0$, l'équation~\eqref{eq:const} devient :
\begin{equation}
  \frac{\hbar^2 k_f^2}{2m^\ast} = \frac{\hbar^2 k_i^2}{2m^\ast} + \hbar\omega_0 , \qquad k_f = k_i + \frac{2\pi}{\lambda_0}.
\end{equation}

La quantité de mouvement du photon est fixée par la longueur d'onde seule : à $\lambda_0 = 1030~\mathrm{nm}$, $k_{\text{ph}} = 2\pi/\lambda_0 \approx 6.1\times10^{6}~\mathrm{m^{-1}}$. La quantité de mouvement de l'électron, en revanche, est fixée par son énergie cinétique dans la bande de conduction. En prenant une valeur représentative $E_i \sim 1~\mathrm{eV}$ (toute énergie de l'ordre de l'eV conduit à la même conclusion), nous obtenons :
\begin{equation}
  k_i = \frac{\sqrt{2 m^\ast E_i}}{\hbar} \approx 5\times10^{9}~\mathrm{m^{-1}}, \qquad \frac{k_{\text{ph}}}{k_i} \approx 10^{-3}.
\end{equation}

La conservation de la quantité de mouvement impose donc $k_f \approx k_i$, d'où $E_f \approx E_i$, ce qui contredit $E_f = E_i + \hbar\omega_0$ sauf si $\hbar\omega_0 \approx 0$.

Le processus devient autorisé si un troisième corps emporte le déséquilibre de quantité de mouvement. Dans un diélectrique polaire comme $\mathrm{SiO_2}$, il s'agit d'un phonon de vecteur d'onde $\vec{q}$ et d'énergie $\hbar\omega_q$, de sorte que :
\begin{equation}
  E_f = E_i + \hbar\omega_0 \pm \hbar\omega_q , \qquad \vec{k}_f = \vec{k}_i + \vec{k}_{\text{ph}} \pm \vec{q}.
\end{equation}

Comme les vecteurs d'onde des phonons couvrent toute la zone de Brillouin :
$$q \lesssim \tfrac\pi a \sim 10^{10}~\mathrm{m^{-1}} \gtrsim k_i$$

La quantité de mouvement peut maintenant être équilibrée, tandis que $\hbar\omega_q \sim 0.06$--$0.15~\mathrm{eV} \ll \hbar\omega_0$ laisse le bilan énergétique essentiellement inchangé. Ce mécanisme assisté par phonons est l'absorption par porteurs libres (FCA), équivalemment le Bremsstrahlung inverse~\cite{mao2004, bulgakova2010}.
